\frontmatter
\pagestyle{plain}

%----------------------------------------------------------------------------------------
%	TITLE PAGE
%----------------------------------------------------------------------------------------

\begin{titlepage}
\begin{singlespacing}
\begin{center}

\vspace*{3mm}
{\LARGE\bfseries \ttitle\par}

\vspace{5mm}
\textit{A thesis submitted in partial fulfillment of the requirements\\for the award of the degree of}

\vspace{4mm}
{\huge\textsc{Master of Technology}}

\vspace{2mm}
\textit{in}

\vspace{2mm}
{\Large\textsc{Information Technology}}

\vspace{3mm}
\textit{with specialization in}

\vspace{2mm}
{\Large\textsc{\specializationname}}

\vspace{5mm}
\safeimage[width=4cm]{Figures/IIITA_Logo.png}{IIITA Logo}
\vspace{5mm}

\textit{Submitted by}\\[2mm]
{\Large\textsc{\studentonename}}\\[1.5mm]
\textit{(\studentoneroll)}

\vspace{4mm}
\textit{Under the Supervision of}\\[2mm]
{\Large\textsc{\supname}}

\vspace{4mm}
\textit{Submitted to}\\[2mm]
{\Large\textsc{\deptname}}

\vspace{6mm}
{\large\textbf{\texthindi{भारतीय सूचना प्रौद्योगिकी संस्थान, इलाहाबाद}}}\\[2mm]
{\large\textsc{Indian Institute of Information Technology, Allahabad}}

\vspace{3mm}
June 2026

\end{center}
\end{singlespacing}
\end{titlepage}

%----------------------------------------------------------------------------------------
%	SUPERVISOR CERTIFICATE
%----------------------------------------------------------------------------------------

\clearpage
\thispagestyle{empty}
\noindent\safeimage[width=\textwidth,keepaspectratio]{Figures/IIITA_Header.png}{IIITA Letterhead}

\vspace{8mm}
\begin{center}
    {\large\bfseries SUPERVISOR CERTIFICATE}
\end{center}
\vspace{6mm}

\noindent It is certified that the work contained in the thesis titled ``\textbf{\ttitle}'', submitted by \textbf{\authorprefix\ \studentonename}, Roll No.\ \textbf{\studentoneroll}, has been carried out under my/our supervision.

\vspace{5mm}
\noindent The work embodied in this thesis has not been submitted elsewhere for the award of any degree.

\vfill

\hfill\begin{minipage}{8cm}
    \rule{\textwidth}{0.4pt}\\[2mm]
    \textbf{\supname}\\
    \deptname\\
    Indian Institute of Information Technology, Allahabad
\end{minipage}

\vspace{15mm}

%----------------------------------------------------------------------------------------
%	CANDIDATE DECLARATION
%----------------------------------------------------------------------------------------

\clearpage
\thispagestyle{empty}
\begin{singlespacing}
\noindent\safeimage[width=\textwidth,keepaspectratio]{Figures/IIITA_Header.png}{IIITA Letterhead}

\vspace{6mm}
\begin{center}
    {\large\bfseries CANDIDATE DECLARATION}
\end{center}
\vspace{4mm}

\noindent I, \textbf{\studentonename}, Roll No.\ \textbf{\studentoneroll}, certify that this thesis titled ``\textbf{\ttitle}'' is submitted by me towards partial fulfillment of the requirements for the award of the degree of \textbf{Master of Technology} in \textbf{Information Technology} with specialization in \textbf{\specializationname}, \textbf{\deptname}, \textbf{Indian Institute of Information Technology, Allahabad}.

\vspace{4mm}
\noindent I understand that plagiarism includes:
\begin{enumerate}
    \item Reproducing someone else's work (fully or partially) or ideas and claiming it as one's own.
    \item Reproducing someone else's work (verbatim copying or paraphrasing) without crediting the original source.
    \item Committing literary theft (copying some unique literary construct).
\end{enumerate}

\noindent I have given due credit to the original authors and sources through proper citation for all the words, ideas, diagrams, graphics, computer programs, experiments, results, and websites that are not my original contributions. I have used quotation marks to identify verbatim sentences and given due credit to the original authors and sources.

\vspace{4mm}
\noindent I affirm that no portion of my work is plagiarized. In the event of a complaint of plagiarism, I shall be fully responsible. I understand that my supervisor may not be in a position to verify that this work is not plagiarized.

\vspace{10mm}
\noindent\begin{minipage}[t]{0.50\textwidth}
    \rule{6cm}{0.4pt}\\[2mm]
    \textbf{\studentonename}\\
    \textbf{\studentoneroll}\\
    \deptname\\
    Indian Institute of Information Technology,\\
    Allahabad\\
    Prayagraj -- 211015, Uttar Pradesh, India
\end{minipage}
\hfill
\begin{minipage}[t]{0.35\textwidth}
    \vspace{2mm}
    Date: \rule{35mm}{0.4pt}
\end{minipage}
\end{singlespacing}

%----------------------------------------------------------------------------------------
%	CERTIFICATE OF APPROVAL
%----------------------------------------------------------------------------------------

\clearpage
\thispagestyle{empty}
\begin{singlespacing}
\noindent\safeimage[width=\textwidth,keepaspectratio]{Figures/IIITA_Header.png}{IIITA Letterhead}

\vspace{4mm}
\begin{center}
    {\large\bfseries CERTIFICATE OF APPROVAL}
\end{center}
\vspace{3mm}

\noindent This is to certify that the thesis entitled ``\textbf{\ttitle}'', submitted by \textbf{\authorprefix\ \studentonename}, Roll No.\ \textbf{\studentoneroll}, in partial fulfillment of the requirements for the award of the degree of \textbf{Master of Technology} in \textbf{Information Technology} with specialization in \textbf{\specializationname} at \textbf{Indian Institute of Information Technology, Allahabad}, has been examined and approved by the Thesis Evaluation Board.

\vspace{3mm}
\noindent The work embodied in this thesis is found to be satisfactory and is approved for the award of the degree of \textbf{Master of Technology}. It is understood that by this approval the undersigned does not necessarily endorse or approve any statement made, opinion expressed, or conclusion drawn therein, but approve the thesis only for the purpose for which it is submitted.

\vspace{4mm}
\noindent Date: \rule{40mm}{0.4pt} \hfill Place: \rule{40mm}{0.4pt}

\vspace{4mm}
\noindent\begin{tabular}{@{}p{0.55\textwidth}p{0.42\textwidth}@{}}
\textbf{\supname} & \\
\textit{Chairperson, Thesis Evaluation Board} & Signature: \rule{38mm}{0.4pt}\\[3.5mm]
\textit{\cosupervisorname} & Signature: \rule{38mm}{0.4pt}\\[3.5mm]
\textbf{\internalexaminername} & \\
\textit{Internal Examiner} & Signature: \rule{38mm}{0.4pt}\\[3.5mm]
\textbf{\externalexaminername} & \\
\textit{External Examiner} & Signature: \rule{38mm}{0.4pt}\\[3.5mm]
\textbf{\hodname} & \\
\textit{Head of Department / Program Coordinator} & Signature: \rule{38mm}{0.4pt}\\[3.5mm]
\textbf{\deanname} & \\
\textit{Dean Academic Affairs} & Signature: \rule{38mm}{0.4pt}\\
\end{tabular}

\vspace{4mm}
\begin{center}
    \textbf{\deptname}\\
    \textit{Indian Institute of Information Technology, Allahabad}\\
    \textit{Prayagraj -- 211015, Uttar Pradesh, India}
\end{center}
\end{singlespacing}

%----------------------------------------------------------------------------------------
%	ABSTRACT PAGE
%----------------------------------------------------------------------------------------

\clearpage
\thispagestyle{plain}
\addchaptertocentry{Abstract}

\begin{center}
    \textsc{Indian Institute of Information Technology, Allahabad}\\[5mm]
    {\large\bfseries \ttitle}\\[5mm]
    \textsc{Abstract}\\[3mm]
    Master of Technology\\
    Department of Information Technology\\[3mm]
    by \studentonename
\end{center}

\vspace{6mm}
Vision-language models (VLMs) endow edge robots with open-vocabulary scene understanding, but their per-inference cost in compute, latency, and energy is prohibitive on resource-constrained accelerators. Mainstream robotic pipelines invoke the VLM on every frame, or on a fixed sampling schedule, regardless of whether the scene has meaningfully changed. Because consecutive frames in a continuous video stream are strongly correlated, this discipline repeatedly re-derives the same understanding of an unchanging scene --- an inefficiency this thesis names \emph{semantic reasoning redundancy}, distinct from the pixel-, frame-, and parameter-level redundancies targeted by prior efficiency research. This thesis proposes a \emph{monitor-then-allocate} architecture in which an always-on, lightweight monitor maintains a composite semantic state from cheap object-detection, multi-object-tracking, and image-embedding signals, and estimates a bounded scalar drift signal $D_t = \alpha D_{\text{embed}} + \beta D_{\text{objects}} + \gamma D_{\text{track}}$ that fuses embedding distance, object-set novelty, and tracking-identity churn against a temporally decayed memory. A three-tier policy maps this signal through two interpretable thresholds to cache reuse (no VLM), lightweight reasoning, or full reasoning, so that expensive multimodal compute is spent only when the scene genuinely changes. The reasoning model is treated as a replaceable black box and the policy is training-free. A modular, edge-oriented reference implementation integrating YOLOv8-n, ByteTrack, CLIP ViT-B/32, and Moondream2 is described, and the trade-off between compute saved and semantic fidelity retained is quantified with a unified \emph{Semantic Compute Efficiency} (SCE) metric against every-frame, uniform-sampling, motion-gating, and embedding-only baselines. In a preliminary evaluation on a continuous stream of $56{,}314$ frames, the proposed three-tier scheduler resolved over $90\%$ of frames from the semantic cache, reducing vision-language model invocations to $9.8\%$ of frames at high semantic retention.

\vspace{5mm}
\noindent\textbf{Keywords:} Vision-language models $\cdot$ Edge robotics $\cdot$ Adaptive compute allocation $\cdot$ Semantic drift $\cdot$ Temporal semantic memory $\cdot$ Selective inference $\cdot$ Object tracking $\cdot$ CLIP embeddings

%----------------------------------------------------------------------------------------
%	ACKNOWLEDGEMENTS
%----------------------------------------------------------------------------------------

\clearpage
\begin{acknowledgements}
\addchaptertocentry{\acknowledgementname}

I express my sincere gratitude to my thesis supervisor, \textbf{\supname}, for their invaluable guidance, constructive feedback, and continuous support throughout the course of this research. Their deep expertise in machine learning and computer vision shaped the technical direction of this work, and their encouragement during challenging phases of experimentation was instrumental in its completion.

\vspace{4mm}
I am grateful to the faculty and staff of the \textbf{\deptname}, \textbf{Indian Institute of Information Technology, Allahabad}, for providing an excellent research environment and computational resources that made this work possible.

\vspace{4mm}
I thank the open-source communities behind PyTorch, Ultralytics YOLO, ByteTrack, CLIP, and Moondream, whose tools and reference implementations enabled the system and the fair, reproducible comparisons reported in this thesis. The broader research community whose foundational work on vision-language modelling, efficient inference, and embodied perception directly enabled the contributions of this thesis deserves acknowledgement.

\vspace{4mm}
Finally, I am deeply grateful to my family and friends for their unwavering support, patience, and encouragement throughout the M.Tech programme.

\end{acknowledgements}

%----------------------------------------------------------------------------------------
%	TABLE OF CONTENTS / FIGURES / TABLES
%----------------------------------------------------------------------------------------

\tableofcontents
\listoffigures
\listoftables

%----------------------------------------------------------------------------------------
%	ABBREVIATIONS
%----------------------------------------------------------------------------------------

\begin{abbreviations}{ll}

\textbf{VLM}   & Vision-Language Model \\
\textbf{VQA}   & Visual Question Answering \\
\textbf{CLIP}  & Contrastive Language-Image Pre-training \\
\textbf{YOLO}  & You Only Look Once (object detector) \\
\textbf{ViT}   & Vision Transformer \\
\textbf{MOT}   & Multi-Object Tracking \\
\textbf{SCE}   & Semantic Compute Efficiency \\
\textbf{FPS}   & Frames Per Second \\
\textbf{GPU}   & Graphics Processing Unit \\
\textbf{CPU}   & Central Processing Unit \\
\textbf{CNN}   & Convolutional Neural Network \\
\textbf{FLOPs} & Floating-Point Operations \\
\textbf{IoU}   & Intersection over Union \\
\textbf{COCO}  & Common Objects in Context (dataset) \\
\textbf{CUDA}  & Compute Unified Device Architecture \\

\end{abbreviations}

%----------------------------------------------------------------------------------------
%	SYMBOLS
%----------------------------------------------------------------------------------------

\begin{symbols}{lll}

$Z_t$               & Composite Semantic State at Frame $t$        & Fused Descriptor \\
$E_t$               & CLIP Image Embedding (unit-norm)             & $\mathbb{R}^{d}$ \\
$O_t$               & Object Class--Count Multiset                 & Symbolic Set \\
$\mathcal{I}_t$     & Active Track Identity Set                    & Identity Set \\
$M_t$               & Temporal Semantic Memory                     & $\mathbb{R}^{d}$ \\
$D_t$               & Composite Semantic Drift                     & $[0,1]$ \\
$D_{\text{embed}}$  & Embedding (cosine) Drift                     & $[0,1]$ \\
$D_{\text{objects}}$& Object-Set Novelty                           & $[0,1]$ \\
$D_{\text{track}}$  & Tracking-Identity Churn                      & $[0,1]$ \\
$\alpha,\beta,\gamma$ & Drift Component Weights ($\sum = 1$)       & Convex Weights \\
$\tau_{\text{low}},\tau_{\text{high}}$ & Tier Decision Thresholds  & $[0,1]$ \\
$\lambda$           & Memory Decay Factor                          & $(0,1)$ \\
$c_t$               & Camera-Motion (Egomotion) Fraction           & $[0,1]$ \\
$\mathrm{SCE}$      & Semantic Compute Efficiency                  & Real Scalar \\

\end{symbols}
